% Chapter 4. Every number is a macro (generated/numbers.tex) or a \needsclaim box.
\chapter{The Cost of Thinking While Listening}
\label{ch:cost}

Speculation on a shared device is paid for by whatever runs next to it. This
chapter prices it: where the silence goes, what co-residence costs before
anything runs, and what each consumer's speculation costs when it is still
running at the endpoint.

\section{Where the silence goes}
\label{sec:stages}

\begin{figure}[t]
  \centering
  \includegraphics[width=\textwidth]{paper_fig1_stage_decomposition.pdf}
  \caption{Median stage durations after the end of speech, reactive pipeline
  without partial decodes, five utterance lengths.}
  \label{fig:th-stages}
\end{figure}

The reactive pipeline, with no partial decodes, starts speaking a median of
\TTFAMulti\,ms after the end of a multi-step problem
(Figure~\ref{fig:th-stages}). The final decode is the largest stage on every
set: \SttShareDev\%, \SttShareShort\%, \SttShareLong\% and \SttShareMulti\% of
the stage sum on the development, short-digit, long-digit and multi-step
sets. The language model's first token takes \LlmShareMulti\% on multi-step
problems.

We had registered that the final's share would exceed half at every length
and rise with length. It does neither at every length: the share is not
monotone in length, and it is below half on two sets. That prediction is
falsified. What stands is that the silence is mostly the recognizer's: its
final decode is the largest single stage everywhere. Speculating the reply
attacks a small stage, and speculating the transcript attacks the largest.

\section{The occupancy tax}
\label{sec:tax}

\begin{figure}[t]
  \centering
  \includegraphics[width=\textwidth]{paper_fig2_occupancy_tax.pdf}
  \caption{Final decode, huge pages kept and huge-page fault fallbacks per
  decode: the recognizer alone and under five co-resident conditions. Six
  runs from one instrumented attribution session, one per condition, 32 turns
  each.}
  \label{fig:th-tax}
\end{figure}

Before any speculation runs, the recognizer is taxed by what shares its
memory. With the recognizer alone, a final decode takes \AloneSTT{}, with
\AloneFallback{} huge-page fault fallbacks per decode, and the process keeps
\AloneHuge\,MB of huge pages (Figure~\ref{fig:th-tax}). Next to it, an inert
process of the language model's size, which holds pages and does nothing
else, raises fallbacks to \BallastFallback{} per decode. The recognizer keeps
\BallastHuge\,MB.

Paired per turn on the same audio, the inert process adds \TaxBallastA{} and
\TaxBallastB{} to the final decode in two attribution sessions. The full
pipeline adds \TaxFullA{} and \TaxFullB. Every interval excludes zero. Each
condition is one run made after its baseline, so time order is confounded
with condition. The two sessions disagree on the size and on the order of
conditions, and only the second was instrumented for huge pages.

The counters point to transparent-huge-page fallback in the decoder's own
allocations. One explanation, offered after the fact and not tested, is that
the decoder allocates and frees working buffers on every call while a
co-resident footprint keeps the pool fragmented. Two per-process mitigations
were tried: loading the recognizer before the language model, and advising
huge pages on its buffers. Neither recovered the huge pages
(\ThpMOne\,MB and \ThpMTwo\,MB of huge pages kept, and \ThpMThree\,MB with both, against \ThpMZero\,MB in the control). The system-wide huge-page policy was left
untouched, so the tax could not be removed with the levers one process
controls; whether a system-level change would remove it is open.

\section{The thinker's overlap penalty}
\label{sec:thinker-penalty}

\begin{figure}[t]
  \centering
  \includegraphics[width=0.85\textwidth]{paper_fig3_overlap_penalties.pdf}
  \caption{Both consumers' overlap penalties on one axis. Thinker bars and the
  partials-removed bar are TTFA; the listener's in-flight bar is final-decode
  time, the fitted cost per second of overlap times the median overlap.}
  \label{fig:th-penalties}
\end{figure}

A reply draft still decoding when the user stops competes with the final
decode and with the reply itself. Against the same item's reactive turns in
the same run, a draft still running at the endpoint delays the first audio by
\ThinkerOverlap{} ($n=\ThinkerOverlapN$ turns). A draft that finished first
shows no detectable penalty in these runs (\ThinkerFinished,
$n=\ThinkerFinishedN$). The penalty depends on device state. Under a
contended memory system it was larger
(\OnsetNinetySixRun{} for a draft still running at budget 96, $n=\OnsetNinetySixRunN$, timed in that test from speech onset). In one contended condition, drafts that had finished before the
endpoint were penalised too (\OnsetSixtyFourFin, $n=\OnsetSixtyFourFinN$). The cost also carries over: speculation in one turn
delays the next
(\Carry{} on the following turn, $n=\CarryN$ turns, two interleaved runs), which is why every comparison since
uses a same-arm washout.

These costs are the thinker's whether or not the draft is useful. Its value is
the subject of Chapter~\ref{ch:value}.

\section{The listener's overlap penalty}
\label{sec:listener-penalty}

The listener pays its own overlap penalty. A partial hypothesis still decoding
at the endpoint runs on the same cores as the final. Across the arms that let
a partial decode run past the endpoint, each second of overlap adds
\ListenerC\,ms of final decode, \ListenerCCI. The interval is taken over
runs, from \PtwoaPop. That lies below the band we registered, 500 to
1500\,ms/s, and is not falsified. The scorer was written after the data. A
coding of the arms chosen afterwards gives a higher value, and is reported
only as a sensitivity.

The cost is visible end to end. Removing the partial decodes from the
reactive pipeline cuts TTFA on the development set by \PartialsDev{}
($n=\PartialsDevN$ items). The canonical reactive pipeline was paying that
for partial transcripts it does not use.

The listener's penalty, unlike the thinker's, can be removed at its source:
the orphaned decode can be stopped. Chapter~\ref{ch:allocator} describes how,
and what the first attempts got wrong.

\section{Residency and memory}
\label{sec:residency}

Keeping the second recognizer resident costs no detectable TTFA: \ResidencyTTFA,
against the same pipeline with the partial engine never built. That headline
drops one repetition run during a breach of the fence. With it, the full pass
gives \ResidencyFull, the same reading. An idle resident engine adds +\ResidencyMem\,MB of agent memory at
run level (range \ResidencyMemRange\,MB over run pairs). The decoding
hypothesis worker, a separate process, holds \WorkerPSS\,MB at spawn.

The residency result rules out one explanation of a cross-date drift in
baseline TTFA, without identifying another. Every claim-bearing comparison
therefore has a same-day baseline.

\section{Energy}
\label{sec:energy-cost}

Speculation costs energy whether or not it pays. Against the reactive
pipeline, continuing a reply draft costs \SpecEnergyr{} per turn, net of the
board's idle floor ($n=\SpecN$ items). The gated listener costs
\EnergyListenerr{} per turn against the reactive pipeline without partial
decodes. Chapter~\ref{ch:allocator} sets that against what it saves.

\section{Summary}

The silence after the user stops is mostly the recognizer's. Its final decode
is taxed by mere co-residence, and either consumer's speculation adds to it if
still running at the endpoint. The thinker's draft costs \ThinkerOverlap{}
when it overlaps the endpoint. The listener's partial costs \ListenerC\,ms of
final decode per second of overlap, but it can be stopped. The next chapter asks
what the thinker's cost buys.
